% draft0: 問題設定 → 研究質問 → 貢献。数値を並べる前に主張の範囲を定める。
\begin{abstract}
    \textbf{draft0：論文設計図。} 各節の箇条書きは執筆メモであり，完成本文ではない．
    概要は「原因サービス順位付けの課題→正常時AR予測とBayes Factor→正式評価→有効範囲と限界」の順にまとめる．
    TODO：正式な手法名・次数・結果表を確定後，本文に即して要約し，仮題を決める．

    \medskip
    \noindent\textbf{論文構成：} 序論／背景／関連研究／提案手法／評価実験／考察／結論．
    付録に補足資料の配置計画，根拠ファイル対応表，次のdraftへのTODOを置く．
\end{abstract}

\clearpage
\pagenumbering{arabic}
% draft0の短い箇条書きがページ境界で分断されないようにする。
\interlinepenalty=10000

\section{序論}
\label{sec:intro}

\subsection{問題設定と研究の動機}
\begin{itemize}
    \item \textbf{主張：} 障害が複数サービスへ波及すると，観測上の大きな変化と障害注入元は一致するとは限らない．本研究は調査候補の順位付けを支援する．
    \item \textbf{書くこと：} メトリクスの時間依存性，静的な前後比較の位置付け，正常時予測からの逸脱を使う動機．図案：原因サービスと影響先を区別した模式図．
    \item \textbf{注意：} 障害と開始時刻は既知，入力はメトリクスとサービス対応，出力はサービス順位．因果グラフの同定や障害検知の達成は主張しない（根拠：\hyperref[ev:protocol]{E2}）．
\end{itemize}

\subsection{研究質問}
\label{sec:rq}
\noindent\textbf{中心RQ：} 正常時のAR予測を基準とした分布変化のBayes Factorは，既知の障害開始時刻の下で，原因サービスの順位付けにどこまで有効か．
\begin{itemize}
    \item \textbf{RQ1（性能）：} 正式RCAEval条件において，既存のメトリクスベース手法と比較した順位精度はどうか．
    \item \textbf{RQ2（構成要素）：} ARの使い方とBFによるスコア化は，Raw入力・GLRTとの比較で何に寄与し，どこで悪化するか．
    \item \textbf{RQ3（有効範囲）：} 単位変換・次数・予測不確実性への感度と，正常時のスコア挙動から，どの限界が分かるか．
\end{itemize}

\subsection{貢献と論文全体の主張}
\label{sec:claims}
\begin{itemize}
    \item \textbf{C1（手法）：} 正常時だけで推定した定常ARの再帰予測誤差を，予測距離に応じて補正し，NIG Bayes Factorでサービス順位へ結び付けるAMBERの構成を示す（\ref{sec:method}節，\hyperref[ev:method]{E1}）．
    \item \textbf{C2（実証）：} 正式評価とケース対応比較によって，構成全体の精度と改善・悪化の内訳を示す（RQ1・RQ2，\ref{sec:experiment}節，\hyperref[ev:main]{E3}--\hyperref[ev:sensitivity]{E6}）．既存手法全般への優越は未確定．
    \item \textbf{C3（限界の明確化）：} 順位付けの有効性とスコアの絶対校正を分け，単位不変性・感度・負の結果から適用範囲を定める（RQ3，\ref{sec:discussion}節，\hyperref[ev:robustness]{E7}・\hyperref[ev:negative]{E8}）．
    \item \textbf{TODO：} 新規性は関連研究照合後に限定して表現する．完成時にC1--C3と結論が一対一で対応するか確認する．
\end{itemize}
